\section{MIDI Domain Logic and Backend Serialization}\label{sec:midi-backend-design}

A primary objective of the implementation is to bridge the gap between high-level musical notation and the low-level
binary encoding required by the MIDI standard.
This section details the mathematical formulas and mapping strategies implemented in the
\texttt{motivo\_music} module to
ensure that the compiler's output is musically accurate.

\subsection{Note to MIDI Pitch Mapping}\label{detail-subsec:note-to-midi-pitch-mapping}

Motivo allows composers to specify pitches using traditional letter names from A to G, accidentals (\#, b), and
octave numbers
(0--8).
The compiler converts these literals into a single numeric byte (0--127) using a deterministic formula.

The implementation defines the \texttt{music::Note} structure as a combination of three fields:

\noindent \textbf{\texttt{Pitch} enum:} Stores the semitone offset within an octave (\texttt{C=0}, \texttt{D=2},
\texttt{E=4}, \texttt{F=5}, \texttt{G=7}, \texttt{A=9}, \texttt{B=11}).

\noindent \textbf{\texttt{Accidental} enum:} Stores a numeric signed offset (\texttt{Flat=-1}, \texttt{Natural=0},
\texttt{Sharp=1}).

\noindent \textbf{\texttt{Octave} integer:} Stores the numeric octave.

The mapping formula used in the \texttt{Note::midi\_number} function is:
\begin{equation}
  P = 12 \times (O + 1) + V_p + V_a\label{eq:midi-pitch-detail}
\end{equation}
Where $P$ is the final MIDI pitch, $O$ is the octave, $V_p$ is the pitch-class value, and $V_a$ is the accidental value.
The constant offset of \textbf{+1} follows the convention used by the compiler and its tests, where \texttt{C4} maps to
MIDI note 60 and \texttt{A4} maps to 69.
By centralizing this calculation at the library level, the compiler ensures that all musical logic throughout the
pipeline is consistent.

\subsection{Instrument and Drum Mapping}\label{detail-subsec:instrument-and-drum-mapping}

The language's \texttt{using} keyword allows the composer to specify the instrumental identity of a track.
This intent is mapped directly to the General MIDI (GM) Level 1 specification.

The implementation utilizes two primary mapping strategies:

\noindent \textbf{Melodic Instruments.} Standard instruments are mapped to their GM program numbers.
For example, \texttt{Piano} maps to 1, \texttt{Guitar} to 24, \texttt{Bass} to 32, and \texttt{Violin} to 40.
These mappings are stored in the \texttt{music::Instrument} enum.
During backend serialization, they are emitted as \textit{Program Change} events on standard MIDI channels (0--8,
11--15).

\noindent \textbf{Percussion Logic.} Unlike melodic tracks, the MIDI standard rigidly reserves \textbf{Channel 10}
specifically for unpitched percussion.
If a track is defined as \texttt{using drums}, the backend automatically routes all its events to Channel 10 and
skips emitting a Program Change event.
Furthermore, individual drum literals in Motivo are mapped to specific GM ``Note-On'' pitches as defined by the
standard drum map:

\noindent \texttt{Kick} $\rightarrow$ 36

\noindent \texttt{Snare} $\rightarrow$ 38

\noindent \texttt{Hihat} $\rightarrow$ 42

\noindent \texttt{Crash} $\rightarrow$ 49

\noindent \texttt{Ride} $\rightarrow$ 51

This explicit mapping aligns drum output with the General MIDI percussion convention in tools that follow that standard.

\begin{table}[htbp]
  \centering
  \renewcommand{\arraystretch}{1.3}
  \begin{tabularx}{\textwidth}{>{\raggedright\arraybackslash}p{3.4cm}X}
    \toprule
    \textbf{Motivo concept} & \textbf{Backend representation} \\
    \midrule
    Tempo declaration & MIDI tempo meta event in track 0. \\
    Time signature declaration & MIDI time-signature meta event in track 0. \\
    Track using piano, guitar, bass, or violin & Separate MIDI track with a Program Change event. \\
    Track using drums & Separate MIDI track routed to the percussion channel. \\
    Note event & Note-On and Note-Off messages with tick positions derived from beats. \\
    Rest & Cursor advancement without emitted note messages. \\
    Chord & Several note events with the same start beat. \\
    Sequence & Sequential flattening of notes, chords, and rests. \\
    \bottomrule
  \end{tabularx}
  \caption{Mapping from Motivo source concepts to MIDI backend artifacts.}\label{tab:motivo-midi-mapping}
\end{table}

\subsection{Binary Serialization of Standard MIDI Files}\label{subsec:midi-binary-serialization}

The final stage of the compilation pipeline translates the flat list of linear \texttt{NoteEvent}s (the IR) into a
Standard MIDI File (SMF) Format 1 binary artifact.
This process requires careful management of binary structures, delta-time calculations, and variable-length encodings.

\subsection{SMF Format 1 Structure}\label{detail-subsec:smf-format-1-structure}

The backend generates a multi-track MIDI file structured into binary chunks.
Each chunk begins with a 4-byte ASCII identifier followed by a 32-bit big-endian integer denoting the chunk's data
length.
The compiler enforces strict big-endian serialization via custom bit-shifting utility functions (\texttt{write\_u16\_be}
and \texttt{write\_u32\_be}).

\noindent \textbf{\texttt{MThd} (Header Chunk):} The file begins with a header specifying \texttt{Format=1}
(multiple simultaneous tracks).
The compiler calculates the total number of tracks as $1 + N$ (where $N$ is the number of instrumental tracks in the
Motivo source) to account for the mandatory tempo track.
The temporal resolution is hardcoded to \textbf{480 PPQ} (Pulses Per Quarter-note), providing a high degree of
precision for complex rhythmic subdivisions.

\noindent \textbf{\texttt{MTrk} (Track Chunks):} The first track (Track 0) is initialized as the ``Tempo Track''.
It contains a \texttt{0xFF 51 03} meta-event.
The backend calculates the \textit{Microseconds Per Quarter-note} ($uspb$) using the formula $uspb =
\frac{60,000,000}{\mathrm{BPM}}$ and writes this as a 24-bit value.
Following the tempo track, each Motivo track is serialized into its own \texttt{MTrk} chunk.

\subsection{The Delta-Time Sorting Algorithm}\label{detail-subsec:the-delta-time-sorting-algorithm}

A fundamental constraint of the MIDI protocol is that it does not utilize absolute timestamps.
Instead, every event within a track must be preceded by a \textit{delta-time}, representing the number of ticks to wait
\textit{after} the immediately preceding event.

Because the IR unrolling process can place events into the track out-of-order (for example, when a \texttt{voice} block
places notes concurrently with the main track via the \texttt{from} keyword), the backend must perform a
re-serialization pass:

\noindent \textbf{Tick Conversion.} Every absolute beat position and duration in the IR is multiplied by the
PPQ constant
to convert it into absolute ticks (\texttt{ticks = beat * 480}).
This creates two tick points for every note: an ``on-tick'' and an ``off-tick''.

\noindent \textbf{Event Separation.} The backend transforms each \texttt{NoteEvent} into two separate
low-level events: a
\texttt{Note-On} event at the start tick, and a \texttt{Note-Off} event at the end tick.

\noindent \textbf{Stable Sorting.} All separated events within a track are collected into a vector and sorted using
\texttt{std::ranges::stable\_sort} by their absolute tick value.
The stable sort preserves the relative order when a \texttt{Note-Off} and a \texttt{Note-On} for the same
pitch occur at the exact
same tick, their relative ordering is preserved, preventing overlapping silence issues.

\noindent \textbf{Delta Computation.} The backend iterates sequentially through the sorted list.
For each event, the delta is computed as the difference between \texttt{current\_tick} and
\texttt{previous\_tick}.

\subsection{Variable-Length Quantities (VLQ)}\label{detail-subsec:variable-length-quantities-(vlq)}

Delta-times in MIDI are compressed using \textit{Variable-Length Quantities}.
This 7-bit encoding scheme allows small delta-times (e.g., zero ticks for simultaneous chord notes) to occupy only one
byte, while larger delays can expand up to four bytes.
The compiler implements a custom \texttt{write\_vlq} function that continuously shifts the 32-bit tick delta by 7 bits,
storing the chunks in a buffer, and applying the high-bit continuation flag (hexadecimal \texttt{80}) to all but the
final byte.

\subsection{Event Emission}\label{detail-subsec:event-emission}

Finally, the backend emits the actual status and data bytes for every event:

\noindent \textbf{Program Change (\texttt{0xC0 | channel}):} Emitted at tick 0 for every melodic track to assign the
chosen instrument.

\noindent \textbf{Note-On (hexadecimal \texttt{90} with the channel encoded in the low nibble):} Contains the pitch and
the defined velocity.

\noindent \textbf{Note-Off (hexadecimal \texttt{80} with the channel encoded in the low nibble):} Explicitly
emitted with
a velocity of 0 to terminate the sound.

\noindent \textbf{End of Track (\texttt{0xFF 2F 00}):} A mandatory meta-event appended to the end of every \texttt{MTrk}
chunk to signal the end of the data stream.

By strictly separating the absolute temporal logic of the IR from the relative delta-time mechanics of the MIDI backend,
the compiler ensures that the musical timeline is consistent with the compiler's beat-to-tick conversion and
that the resulting binary file targets the Standard MIDI File structure.
